\documentclass[10pt,a4paper]{amsart}
\usepackage[margin=19mm]{geometry}
\usepackage{amsmath,amssymb,mathtools}
\usepackage[scr=rsfs,cal=euler]{mathalfa}
\usepackage{xfrac}
\usepackage[unicode,hidelinks,bookmarks=false]{hyperref}
\newcommand{\ie}{i.e.\ }
\newcommand{\cf}{cf.\ }
\newcommand{\resp}{respectively\ }
\newcommand{\Subsection}{\subsection}
\providecommand{\nobreakdash}{\nobreak\hbox{-}}
\setlength{\emergencystretch}{2em}
\multlinegap0pt
\usepackage{bm}
\usepackage{bbm}
\usepackage{dsfont}
\usepackage{xspace}
\usepackage{cancel}
\usepackage{mathtools}
\usepackage[shortlabels]{enumitem}
\usepackage{amsthm}
\makeatletter
\@ifundefined{lem}{\newtheorem{lem}{Lemma}}{}
\makeatother
\title[Quasi-strongly regular graphs on the flags of symmetric desi]{Quasi-strongly regular graphs on the flags of symmetric designs}
\author{Eugenia O'Reilly-Regueiro, Octavio B. Zapata-Fonseca}
\date{Source: arXiv:2505.21272v1 (CC BY 4.0)}
\begin{document}
\maketitle

\noindent\emph{Source and licence.} Quasi-strongly regular graphs on the flags of symmetric designs. Version arXiv:2505.21272v1 (CC BY 4.0), distributed under the Creative Commons Attribution 4.0 International licence, as recorded in the arXiv licence metadata for this exact version and shown on the abstract page https://arxiv.org/abs/2505.21272v1. Full rights notice: licenses/arxiv-2505.21272-CC-BY-4.0.txt.\par\smallskip

\noindent\textit{Editorial scope.} Counted unit: the Lem whose source label is \texttt{beta} in the article (source0.tex lines 254--262, source label \texttt{beta}), delivered together with the article's own complete original proof of it (source0.tex lines 264--313, one \texttt{proof} environment of 50 lines). No mathematics is written, repaired or re-derived: statement and proof are the article's own text line for line. Required context retained uncounted: none. Every definition, display and label of those lines is retained unchanged. Omitted from this package and not redistributed: 1--253, 263--263, 314--496. Nothing inside a retained excerpt is omitted or altered: every retained body is delivered exactly as the article's lines are written, so no omission receipt of any kind accompanies this package. Citations: the retained excerpts carry no citation command at all, so no bibliography entry of the article is transcribed and no reference is redistributed. Reference adaptation: none was needed and none was made; every reference inside the delivered unit resolves inside it, so no unresolved reference or citation is produced. Printed slips kept literal: no printed word, symbol, hypothesis or number of the counted statement or of its proof was altered. Layout and class: the article's own class is \texttt{article} with packages including amsmath, amssymb, amsthm, theoremref, babel, xcolor, hyphenat, hyperref, enumerate; the wrapper is the lane's pinned \texttt{amsart} editorial wrapper at 10pt on A4, which restates the theorem-like environments the retained text needs and adds the article's own macro definitions that the retained text needs (none), with extra wrapper packages (bm, bbm, dsfont, xspace, cancel, mathtools, enumitem). The delivered numbering is the unit's own. Assets: no retained span contains a figure environment, a graphics-inclusion command, a PDF-page inclusion, a table or a TikZ picture; no image file is copied into this package, the package declares no asset dependency and no image file is redistributed with it. Licence: version arXiv:2505.21272v1 of this article is distributed under the Creative Commons Attribution 4.0 International licence, as recorded in the arXiv licence metadata for this exact version and shown on the abstract page https://arxiv.org/abs/2505.21272v1. A full rights notice is delivered at licenses/arxiv-2505.21272-CC-BY-4.0.txt. Characters: every retained excerpt is pure ASCII, so the delivered engine, XeLaTeX with its default Latin Modern text font, reports no missing character.
\begin{lem}~\label{beta}
    Let $D=(P,\mathcal{B})$ and $D'=(P',\mathcal{B}')$ be two non-trivial $(v,k,2)$-biplanes, and $\Gamma_2(D)$ and $\Gamma_2(D')$ the corresponding graphs on their flags, and suppose $\Gamma_2(D)\cong\Gamma_2(D')$. If $\beta:V(\Gamma_2(D))\to V(\Gamma_2(D'))$ is such an isomorphism, then $\beta$ induces a bijection $\beta_P:P\to P'$ and a bijection $\beta_{\mathcal{B}}:\mathcal{B}\to\mathcal{B}'$ such that
     for any $p_1,p_i,p_j\in P$ and $c_1,c_i,c_j\in\mathcal{B}$ such that $(p_1,c_i),(p_1,c_j),(p_i,c_1),(p_j,c_1)\in V(\Gamma_2(D))$, the following hold:
    \begin{enumerate}[(i)]
        \item~\label{b1} If $\beta(p_1,c_i)=(p_1',c_i')$ and $\beta(p_1,c_j)=(q',c_j')$, then $p_1'=q'$, and
        \item~\label{b2} If $\beta(p_i,c_1)=(p_i',c_1')$ and $\beta(p_j,c_1)=(p_j',d')$, then $c_1'=d'$,
    \end{enumerate}
    that is, for any $p\in P,c\in\mathcal{B}$, if $p\in c$ then $\beta(p,c)=(\beta_P(p),\beta_\mathcal{B}(c))$.
\end{lem}

\begin{proof}
    Let $D=(P,\mathcal{B})$ and $D'=(P',\mathcal{B}')$ be two non-trivial $(v,k,2)$-biplanes, and $\Gamma_2(D)$, $\Gamma_2(D')$ the corresponding graphs on their flags, and suppose there is an isomorphism $\beta:V(\Gamma_2(D))\to V(\Gamma_2(D'))$, such that for any $(p,c),(q,d)\in V(\Gamma_2(D))$, $(p,c)\sim(q,d)\iff\beta(p,c)\sim\beta(q,d)$.\\
    
    For~(\ref{b1}), let $c_1\in\mathcal{B}$ with $p_1,p_i,p_j\in c_1$ such that $\{p_1,p_i\}=c_1\cap c_i$ and $\{p_1,p_j\}=c_1\cap c_j$ for some unique $c_i,c_j\in\mathcal{B}$, and such that $\beta(p_1,c_i)=(p_1',c_i')$ and $\beta(p_1,c_j)=(q',c_j')$. Then there is a point $x_{ij}\in P$ such that $c_i\cap c_j=\{p_1,x_{ij}\}$. This point $x_{ij}$ might or might not be in $c_1$, however we do have the following adjacencies in $\Gamma_2(D)$:\\

$(p_1,c_i)\sim(x_{ij},c_j)\iff\beta(p_1,c_i)\sim\beta(x_{ij},c_j)\iff(p_1',c_i')\sim\beta(x_{ij},c_j).$\\
    
    Also:\\
    
$(p_1,c_j)\sim(x_{ij},c_i)\iff\beta(p_1,c_j)\sim\beta(x_{ij},c_i)\iff(q',c_j')\sim\beta(x_{ij},c_i)$.\\

    If $p_1\neq q'$, then there are two blocks $d_1',d_2'\in\mathcal{B}'$ such that $d_1'\cap d_2'=\{p_1',q'\}$, and there are two adjacent pairs of vertices in $\Gamma_2(D')$: $(p_1',d_1')\sim(q',d_2')$ and $(p_1',d_2')\sim(q',d_1')$.\\

    Since $\beta$ is an isomorphism, these vertices must correspond to two pairs of adjacent vertices in $\Gamma_2(D)$ given by the elements of $(p_1,c_i)$ and $(p_1,c_j)$. Since $p_1$ is the same point of $D$ in both vertices, they must then correspond to the intersection of $c_i\cap c_j=\{p_1,x_{ij}\}$. That is:
    \begin{align*}
    &\{\beta^{-1}(p_1',d_1'),\beta^{-1}(p_1',d_2'),\beta^{-1}(q',d_1'),\beta^{-1}(q',d_2')\}\\
    &\qquad\qquad=\{(p_1,c_i),(p_1,c_j),(x_{ij},c_i),(x_{ij},c_j)\}.
\end{align*}
    Without loss of generality, suppose $c_i'=d_1'$ and $c_j'=d_2'$. We have the following:\\

$(p_1,c_i)\sim(x_{ij},c_j)\iff\beta(p_1,c_i)\sim\beta(x_{ij},c_j)\iff(p_1',c_i')\sim\beta(x_{ij},c_j)\iff(p_1',d_1')\sim(q',d_2')$, so $\beta(x_{ij},c_j)=(q',d_2')=(q',c_j')=\beta(p_1,c_j)$, which is a contradiction.\\

The proof of~(\ref{b2}) is very similar. Now suppose there is a block $c_1\in\mathcal{B}$ with $p_i,p_j\in c_1$ such that $\beta(p_i,c_1)=(p_i',c_1')$ and $\beta(p_j,c_1)=(p_j',d')$ with $c_1'\neq d'$.\\

Since $k>3$, as before, let $p_1\in c_1$ be another point, and $c_i,c_j\in\mathcal{B}$ be the blocks such that $\{p_1,p_i\}=c_1\cap c_i$ and $\{p_1,p_j\}=c_1\cap c_j$.\\

Again, we have the following adjacencies in $\Gamma_2(D)$:\\

$(p_1,c_i)\sim(p_i,c_1),$ and $(p_1,c_j)\sim(p_j,c_1)$.\\

This happens if and only if:\\

$\beta(p_1,c_i)\sim\beta(p_i,c_1)=(p_i',c_1'),$ and $\beta(p_1,c_j)\sim\beta(p_j,c_1)=(p_j',d')$.\\

Since $c_1'\neq d'$, there are two points $q_1',q_2'\in P'$ such that $c_1'\cap d'=\{q_1',q_2'\}$. This induces two pairs of adjacent vertices in $\Gamma_2(D')$:\\

$(q_1',c_1')\sim(q_2',d')$, and $(q_1',d')\sim(q_2',c_1')$.\\

Since $\beta$ is an isomorphism, these must correspond to two pairs of adjacent vertices in $\Gamma_2(D)$ induced by the elements of the flags $(p_i,c_1)$ and $(p_j,c_1)$. Since the block $c_1$ is the same in both vertices, they must correspond to the intersection $\{p_i,p_j\}=c_1\cap d$ for some block $d\in\mathcal{B}$. Therefore:\\

$\{\beta(p_1,c_i),\beta(p_1,c_j),\beta(p_i,c_1),\beta(p_j,c_1)\}=\{(q_1',c_1'),(q_2',c_1'),(q_1',d'),(q_2',d')\}$.\\

Suppose without loss of generality that $p_i'=q_1'$. Then:\\

$\beta(p_1,c_i)\sim\beta(p_i,c_1)=(p_i',c_1')=(q_1',c_1')\sim(q_2',d')$, so $\beta(p_1,c_i)=(q_2',d')$.\\

From above:\\

$\beta(p_1,c_j)\sim\beta(p_j,c_1)=(p_j',d')=(q_2',d')$, a contradiction which completes the proof.
\end{proof}

\end{document}
